\section*{NeurIPS Paper Checklist}

\begin{enumerate}

\item {\bf Claims}
    \item[] Question: Do the main claims made in the abstract and introduction accurately reflect the paper's contributions and scope?
    \item[] Answer: \answerYes{}
    \item[] Justification: The abstract and Introduction state the central claim---that confidence fails to fall when visual evidence is removed, replaced, or corrupted, so calibration error rises---and the Results (RQ1--RQ5) and Tables~\ref{tab:main}--\ref{tab:verbalized} support exactly these claims, with scope limited to the three evaluated models and two datasets.

\item {\bf Limitations}
    \item[] Question: Does the paper discuss the limitations of the work performed by the authors?
    \item[] Answer: \answerYes{}
    \item[] Justification: The Discussion contains a dedicated Limitations paragraph covering the three-model/two-dataset scope, greedy decoding, the mean-over-tokens confidence proxy, and small-split ECE noise (for which a bootstrap interval is reported).

\item {\bf Theory assumptions and proofs}
    \item[] Question: For each theoretical result, does the paper provide the full set of assumptions and a complete (and correct) proof?
    \item[] Answer: \answerNA{}
    \item[] Justification: The paper is an empirical audit and contains no theoretical results or proofs.

    \item {\bf Experimental result reproducibility}
    \item[] Question: Does the paper fully disclose all the information needed to reproduce the main experimental results of the paper to the extent that it affects the main claims and/or conclusions of the paper (regardless of whether the code and data are provided or not)?
    \item[] Answer: \answerYes{}
    \item[] Justification: The Setup describes the exact models, datasets, image conditions, decoding settings (greedy, at most 32 new tokens), confidence definition, and ECE binning; the evaluation harness is released, and all datasets and model checkpoints are public.

\item {\bf Open access to data and code}
    \item[] Question: Does the paper provide open access to the data and code, with sufficient instructions to faithfully reproduce the main experimental results, as described in supplemental material?
    \item[] Answer: \answerYes{}
    \item[] Justification: We release the evaluation harness (model-agnostic inference, scoring, stratification, bootstrap CIs, and plotting scripts) with usage instructions; both datasets (VQA-RAD, PathVQA) and all three model checkpoints are publicly available on the Hugging Face Hub as identified in the bibliography.

\item {\bf Experimental setting/details}
    \item[] Question: Does the paper specify all the training and test details (e.g., data splits, hyperparameters, how they were chosen, type of optimizer) necessary to understand the results?
    \item[] Answer: \answerYes{}
    \item[] Justification: No training is performed; the Setup specifies the test splits (VQA-RAD $n{=}451$, PathVQA $n{=}6719$), the three image conditions, greedy decoding with at most 32 new tokens, and the 10-bin ECE used throughout.

\item {\bf Experiment statistical significance}
    \item[] Question: Does the paper report error bars suitably and correctly defined or other appropriate information about the statistical significance of the experiments?
    \item[] Answer: \answerYes{}
    \item[] Justification: For the key RQ3 result we report a bootstrap 95\% confidence interval (5000 resamples of items with replacement, percentile interval) on the intact-to-mismatch ECE increase for closed questions; the interval excludes zero. Other metrics are computed on the full test splits.

\item {\bf Experiments compute resources}
    \item[] Question: For each experiment, does the paper provide sufficient information on the computer resources (type of compute workers, memory, time of execution) needed to reproduce the experiments?
    \item[] Answer: \answerYes{}
    \item[] Justification: All experiments are inference only (no training). Each model/dataset run used a single NVIDIA A100 GPU; the heaviest runs (the full severity sweep over PathVQA, seven conditions $\times$ 6719 items) take on the order of hours, and total compute is modest.

\item {\bf Code of ethics}
    \item[] Question: Does the research conducted in the paper conform, in every respect, with the NeurIPS Code of Ethics \url{https://neurips.cc/public/EthicsGuidelines}?
    \item[] Answer: \answerYes{}
    \item[] Justification: The work uses publicly available, de-identified medical VQA datasets and performs inference only; it involves no human subjects and is aimed at improving the safety and reliability of medical VLMs.

\item {\bf Broader impacts}
    \item[] Question: Does the paper discuss both potential positive societal impacts and negative societal impacts of the work performed?
    \item[] Answer: \answerYes{}
    \item[] Justification: The Discussion (Clinical implications) describes the positive impact---surfacing a safety-relevant calibration failure invisible to accuracy-only benchmarks and recommending cheap safeguards---and the motivating negative scenario of confidently wrong answers on degraded or mismatched medical images.

\item {\bf Safeguards}
    \item[] Question: Does the paper describe safeguards that have been put in place for responsible release of data or models that have a high risk for misuse (e.g., pre-trained language models, image generators, or scraped datasets)?
    \item[] Answer: \answerNA{}
    \item[] Justification: We release only an evaluation harness (analysis code); no new model, generator, or scraped dataset with misuse risk is released.

\item {\bf Licenses for existing assets}
    \item[] Question: Are the creators or original owners of assets (e.g., code, data, models), used in the paper, properly credited and are the license and terms of use explicitly mentioned and properly respected?
    \item[] Answer: \answerYes{}
    \item[] Justification: All datasets and models are cited, and their licenses/terms are noted in the bibliography (VQA-RAD CC0-1.0, PathVQA MIT, Qwen2.5-VL and LLaVA-1.6 Apache-2.0, MedGemma under the Health AI Developer Foundations terms).

\item {\bf New assets}
    \item[] Question: Are new assets introduced in the paper well documented and is the documentation provided alongside the assets?
    \item[] Answer: \answerYes{}
    \item[] Justification: The only new asset is the evaluation harness, which is released with a README describing dependencies, the datasets/models to download, and exact commands to reproduce each experiment.

\item {\bf Crowdsourcing and research with human subjects}
    \item[] Question: For crowdsourcing experiments and research with human subjects, does the paper include the full text of instructions given to participants and screenshots, if applicable, as well as details about compensation (if any)?
    \item[] Answer: \answerNA{}
    \item[] Justification: The paper does not involve crowdsourcing or research with human subjects; it uses existing public datasets.

\item {\bf Institutional review board (IRB) approvals or equivalent for research with human subjects}
    \item[] Question: Does the paper describe potential risks incurred by study participants, whether such risks were disclosed to the subjects, and whether Institutional Review Board (IRB) approvals (or an equivalent approval/review based on the requirements of your country or institution) were obtained?
    \item[] Answer: \answerNA{}
    \item[] Justification: The paper does not involve human subjects, so no IRB approval is required.

\item {\bf Declaration of LLM usage}
    \item[] Question: Does the paper describe the usage of LLMs if it is an important, original, or non-standard component of the core methods in this research? Note that if the LLM is used only for writing, editing, or formatting purposes and does \emph{not} impact the core methodology, scientific rigor, or originality of the research, declaration is not required.
    \item[] Answer: \answerNA{}
    \item[] Justification: The vision-language models in this paper are the objects of evaluation, not tools used to develop the methodology; no LLM is an important, original, or non-standard component of the core method.

\end{enumerate}
